\documentclass{article}
\usepackage{graphicx} % Required for inserting images
\usepackage{array}   % Para mejor control de columnas
\usepackage{booktabs} % Para líneas horizontales más profesionales (opcional)

\title{MIPS}
\author{Victor Pinto}
\date{May 2026}

\begin{document}

\maketitle

\section{Que es Mips 32?}
Mips 32 es una arquitectura de 32 Bits para microprocesadores que sigue la filosofia RISC (conjunto de instrucciones simples o reducidas)

\section{Primeros ladrillos}
Antes de adentrarnos al como esta organizada la estructura de Mips, hay que entender algunos conceptos basicos:

    Palabra = Unidad de informacion, equivalente a 4 Bytes \\
    Registro = Bloque de acceso rapido que tiene Mips para datos \\
    Memoria = Espacio en donde Mips puede almacenar o sustraer datos \\
    Alineamiento = Mips trabaja un alineamiento de 4 Bytes, es decir, se mueve siempre en 4 unidades \\
    
\section{Registros}
Mips cuenta con 32 Registros de 4 Bytes cada uno para realizar operaciones de manera rapida. Dentro de su arquitectura, los registros se encargan de diferentes tareas segun el tipo de registro \\
En los registros podemos cargar la informacion presente en la memoria en estos, y utilizarlos directamente en las instrucciones. De igual forma, los registros pueden almacenar informacion para llevarla la memoria \\
Existen registros cuyo proposito es de almacenar informacion de manera temporal, de guardar direcciones de retorno, o de emplearse para el paso de parametros. Asi, podemos ver como cada registro cumple un papel durante la ejecucion de un programa ensamblado, mas alla de solo almacenar informacion de rapido alcance para el procesador \\
A continuacion se anexa una tabla en donde se encuentra el numero de cada registro, el nombre y una breve descripcion sobre el funcionamiento que este presenta \\ \\

\begin{table}[h]
\centering
\begin{tabular}{|c|l|p{7cm}|}
\hline
\textbf{Numero} & \textbf{Nombre} & \textbf{Descripcion} \\ \hline
0     & \$zero  & Almacena siempre la constante 0 \\ \hline
1     & \$at    & Reservado para el sistema \\ \hline
2-3   & \$v0-\$v1 & Valores de retorno y de expresion \\ \hline
4-7   & \$a0-\$a3 & Registros para argumentos o parametros \\ \hline
8-15  & \$t0-\$t7 & Registros temporales de uso general \\ \hline
16-23 & \$s0-\$s7 & Registros temporales que pueden ser salvados \\ \hline
24-25 & \$t8-\$t9 & Temporales adicionales \\ \hline
26-27 & \$k0-\$k1 & Registros Kernel, reservados para el sistema \\ \hline
28    & \$gp    & Puntero Global al segmento de datos estaticos \\ \hline
29    & \$sp    & Puntero de Pila, a la ultima posicion de si \\ \hline
30    & \$fp    & Puntero de marco \\ \hline
31    & \$ra    & Direccion de retorno, almacena la direccion a la que se debe retornar \\ \hline
\end{tabular}
\end{table}

Un Registro especial: El Registro PC \\
Se trata de un registro llamado Program Counter, el cual almacena siempre la direccion de memoria de la instruccion que se este ejecutando en el momento \\

\section{Instrucciones}
Una instruccion es una palabra presente en el hardware que permite efectuar una operacion basica, generalmente con pocos operandos. Las instrucciones son puntuales y especificas, lo que permite un control sobre lo que sucede al nivel de la maquina.\\
Las instrucciones son clasificadas en formatos, que representan su estructuracion interna y la tarea que pueden cumplir.
Los formatos son: \\

    \textbf{I - Format} = Usado para instrucciones de transferencia de datos inmediatos

    \textbf{R - Format} = Usado para el trabajo con registros 

    \textbf{J - Format }= Se usa para saltar a otras posiciones \\
    
A continuacion se anexa una tabla con instrucciones basicas:

\begin{table}[h]
\centering
\begin{tabular}{|c|c|p{6.5cm}|c|}
\hline
\textbf{Nombre} & \textbf{Formato} & \textbf{Descripcion} & \textbf{Ejemplo} \\ \hline
add  & R & Se almacena en el primer registro la suma de los valores de los otros dos operandos & add \$s1,\$s2,\$s3 \\ \hline
addi & I & Se almacena en el primer registro la suma de otro registro con una constante & addi \$s1,\$s2,20 \\ \hline
lw   & I & Carga una palabra de la memoria y la lleva a un registro, especificando un desplazamiento & lw \$t1,100(\$s2) \\ \hline
sw   & I & Almacena una palabra de un registro a la memoria, especificando desplazamiento & sw \$t1,200(\$s2) \\ \hline
beq  & I & Salta a una etiqueta X si dos registros son iguales & beq \$t1,\$t2,X \\ \hline
bne  & I & Salta a una etiqueta si dos registros no son iguales & bne \$s1,\$s2,X \\ \hline
slt  & R & Fija en 1 el primer registro si el segundo es menor que el tercero, 0 en caso contrario & slt \$t0,\$s1,\$s2 \\ \hline
slti & I & Al igual que con slt, compara el registro si el registro es menor pero con una constante & slti \$t0,\$s1,100 \\ \hline
j    & J & Realiza un salto incondicional a una direccion especifica & j 5000 \\ \hline
sll  & R & Desplaza bits cierta cantidad de veces segun una cantidad a la izquierda & sll \$s1,\$s2,10 \\ \hline
\end{tabular}
\end{table}

\end{document}





